\documentclass[12pt]{article}
\usepackage{polyglossia}
\usepackage{fontspec}
\usepackage[export]{adjustbox}
\usepackage{amsfonts, amsmath, amssymb}
\usepackage{annotate-equations}
\usepackage{bm}
\usepackage{bookmark}
% \usepackage{breqn}
\usepackage{derivative}
\usepackage{empheq}
\usepackage[inline]{enumitem}
\usepackage{fancyhdr}
\usepackage{geometry}
\usepackage{graphicx}
\usepackage{kantlipsum}
\usepackage{lastpage}
\usepackage{linebreaker} % line-breaker algorithm in LuaLaTeX
\usepackage{lualatex-math} % Fixes for mathematics
\usepackage{mathtools}
\usepackage{microtype}
\usepackage{mismath}
\usepackage{nicefrac}
\usepackage[hyperref]{ntheorem}
% \usepackage{phfqit} % Bra-ket notation
% \usepackage{physics}
\usepackage[final]{showlabels}
\usepackage{simples-matrices}
\usepackage{siunitx}
\usepackage{soul}
\usepackage{spbmark}
\usepackage{titling}
\usepackage{tabularray}
\usepackage{tasks}
\usepackage{tensor}
\usepackage{xcolor}
\usepackage{xfrac}
\usepackage{xparse}
\usepackage{xspace}
\usepackage{hyperref}
% Polyglossia config
\setdefaultlanguage[variant=mexican]{spanish}
\setotherlanguage{english}
\usepackage{cleveref}

%%%%%%%%%%%%%%%%%%%%%%%%%%%%%%%%%%%%%%%%%%%%%%%%%%%%%%%%%%
%%%%%%%%%%%%%%%%%%%%%%%%%%%%%%%%%%%%%%%%%%%%%%%%%%%%%%%%%%
%%%%%%%%%%%%%%%%%%Configuraciones extras%%%%%%%%%%%%%%%%%%
\definecolor{base3}{RGB}{253, 246, 227}%
\definecolor{pinkwave}{RGB}{255, 0, 128}%
\definecolor{customBlue}{RGB}{11, 61, 98}%
\pagecolor{base3}
\graphicspath{{img/}}
\setlength{\parindent}{2em} % Sangría
\setlength{\parskip}{0.5em} % Espacio entre párrafos
\linespread{1.1} % line spacing
\setlength{\jot}{10pt} % Space between lines in multiline eqs
% \crefname{equation}{ec.}{ecs.} % Equation's cross-reference name
\crefname{equation}{}{} % Equation's cross-reference name

% Line-breaker config
\linebreakersetup {
    maxtolerance = 90,
    maxemergencystretch = 1em,
    maxcycles = 4
}

%%%%%%%%%%%%%%%%%%Theorem environments%%%%%%%%%%%%%%%%%%
% Configuración de ambiente para problema
\theoremstyle{break}
\theoremheaderfont{\Large\normalfont\bfseries}
\theorembodyfont{\normalfont}
\theoremseparator{\bigskip} % Spacing between header and body
\theorempreskip{1.5em}
\theorempostskip{\topsep\bigskip}
\theorempostwork{
    \color{customBlue} \hrule width \hsize height 2pt \kern 1mm \hrule width \hsize
    }
\newtheorem{exercise}{Problema}
\counterwithin{equation}{exercise}

\makeatletter
\LetLtxMacro\old@exercise\exercise
\renewcommand*{\exercise}{%
  \@ifnextchar[{%
    \exercise@opt
  }{%
    \old@exercise
    \exercise@bookmark{}%
  }%
}
\def\exercise@opt[#1]{%
  \old@exercise[{#1}]%
  \exercise@bookmark{: #1}%
}
\newcommand*{\exercise@bookmark}[1]{%
  \bookmark[dest=\@currentHref]{Problema \theexercise#1}%
}
\makeatother

% Configuración de ambiente para solución
\theoremstyle{nonumberbreak}
\theoremheaderfont{\Large\normalfont\bfseries}
\theorembodyfont{\normalfont}
\theoremseparator{\medskip}
\theorempreskip{1em}
\theorempostskip{\topsep\medskip}
\newtheorem{solution}{Solución}

%%%%%%%%%%%%%%%%%%%%%%%%%%%%%%%%%%%%%%%%%%%%%%%%%%%%%%%%


% Configuración del paquete hyperref
\hypersetup{
    colorlinks = true,%
    linkcolor={[rgb]{0,0.2,0.6}},%
    citecolor={[rgb]{0,0.6,0.2}},%
    filecolor={[rgb]{0.8,0,0.8}},%
    urlcolor={[rgb]{0.8,0,0.8}},%
    runcolor={[rgb]{0.8,0,0.8}},% 
    menucolor={[rgb]{0,0.2,0.6}},%
    linkbordercolor={[rgb]{0,0.2,0.6}},%
    citebordercolor={[rgb]{0,0.6,0.2}},%
    filebordercolor={[rgb]{0.8,0,0.8}},%
    urlbordercolor={[rgb]{0.8,0,0.8}},%
    runbordercolor={[rgb]{0.8,0,0.8}},%
    menubordercolor={[rgb]{0,0.2,0.6}},% 
    pdftitle={Tarea 12},%
    pdfauthor={López Merino Marcos},%
    pdfsubject={Relatividad},%
    pdfkeywords={Facultad de Ciencias, UNAM, Relatividad},%
    unicode = true%
}

%%%%%%%%%%%%%%%%%%siunitx configuration%%%%%%%%%%%%%%%%%%
% Configuración del paquete siunitx
\sisetup{
	output-decimal-marker = {.}, 
	per-mode = symbol-or-fraction,
	separate-uncertainty = false,
	exponent-product = \cross,
    % inter-unit-product = \ensuremath{{}\vdot{}}
}

% Declaring new units
\DeclareSIUnit\kilogram{\kilo\gram}

%%%%%%%%%%%%%%%%%%%%%%%%%%%%%%%%%%%%%%%%%%%%%%%%%%%%%%%%

% Geometría del documento
\geometry{
    letterpaper,
    top = 0.6in,
    bottom = 0.8in,
    left = 0.6in,
    right = 0.6in,
    footskip = 38pt
}

%%%%%%%%%%%%%%%%%%Nuevos comandos%%%%%%%%%%%%%%%%%%
\newcommand*{\group}{8231}
\newcommand*{\classname}{Relatividad}
\newcommand*{\homeworknumber}{Tarea 12}
\newcommand*{\name}{Marcos López Merino}

% unit vector i
\newcommand*{\uveci}{{\bm{\hat{\textnormal{\bfseries\i}}}}}
% unit vector j
\newcommand*{\uvecj}{{\bm{\hat{\textnormal{\bfseries\j}}}}}
% unit vector
\DeclareRobustCommand{\uvec}[1]{{%
  \ifcsname uvec#1\endcsname
     \csname uvec#1\endcsname
   \else
    \bm{\hat{\mathbf{#1}}}%
   \fi
}}% 
\newcommand{\idest}{\emph{i.e.},\xspace} % id est
% Espacio vectorial, e.g., ℝ, ℂ, ℕ, etc.
\NewDocumentCommand{\vecspace}{m o}{%
  \IfValueTF{#2}{%
    \mathbb{#1}^{#2}%
  }{%
    \mathbb{#1}%
  }%
}

\newcommand*{\observer}{\mathcal{O}}
\newcommand*{\primeobserver}{\mathcal{O}^{\prime}}
\newcommand*{\biprimeobserver}{\mathcal{O}^{\prime\prime}}
\newcommand*{\triprimeobserver}{\mathcal{O}^{\prime\prime\prime}}
\newcommand*{\inlinesol}{\vspace*{10pt}\textbf{Solución}\vspace*{10pt}}
\renewcommand*{\grad}{\nabla}
\renewcommand*{\div}{\nabla \cdot}
\renewcommand*{\curl}{\nabla \mul}
\newcommand*{\laplacian}{\nabla^{2}}
\newcommand*{\qq}{\quad\quad}
\DeclareDifferential{\cdif}{\grad}[style-notation=single]

% \newcommand*{\change}[1]{\Delta{#1}\xspace}

\NewDocumentCommand{\change}{ s o m }{%
    \IfBooleanTF{#1}{%
        \Delta {#3^{\prime}}
    }{%
        \Delta {#3}
    }
}
\NewDocumentCommand{\interval}{ s o m }{%
    \IfBooleanTF{#1}{%
        (\Delta {#3^{\prime}})^{2}
    }{%
        (\Delta {#3})^{2}
    }
}
\DeclarePairedDelimiterX\set[1]\lbrace\rbrace{#1}
% \newcommand{\crefrangeconjunction}{--}
% \newcommand{\crefpairconjunction}{\xspace y\xspace}

% Defining a variant of Aboxed command from mathtools
\makeatletter
\newcommand*\Acolorboxed[2][pinkwave]{%
   \let\bgroup{\romannumeral-`}%
   \@Acolorboxed{#1}#2&&\ENDDNE
}
\def\@Acolorboxed#1#2&#3&#4\ENDDNE{%
  \ifnum0=`{}\fi
  \setbox\z@\hbox{$\displaystyle#2{}\m@th$\kern\fboxsep \kern\fboxrule}%
  \edef\@tempa{\kern\wd\z@ & \kern-\the\wd\z@ \fboxsep\the\fboxsep \fboxrule\the\fboxrule}%
  \@tempa
  \fcolorbox{#1}{base3}{\m@th$\displaystyle#2#3$}%
} 
\makeatother

%%%%%%%%%%%%%%%%%%Portada y configuración%%%%%%%%%%%%%%%%%%
% Configuración de portada
\setlength{\droptitle}{-60pt} % raise the title

% Portada
\title{
    \textbf{\homeworknumber}\\
    \normalsize\vspace{0.1in}\small{\textbf{Entrega}:~\today}
    \vspace{-1.5in}
}
\author{}
\date{}

%%%%%%%%%%%%%%%%%%Header and footer%%%%%%%%%%%%%%%%%%
\setlength{\headheight}{15.2pt}
\pagestyle{fancy}
\lhead{Grupo \group, Sem. 2023-1}
\chead{\classname}
\rhead{\name}
\lfoot{\includegraphics[scale = 0.2, valign = c]{LogoFCUNAMcolor.pdf}}
% \lfoot{\includegraphics[scale = 0.1, valign = c]{example-image}}
\cfoot{\homeworknumber}
\rfoot{Pág. \thepage \hspace{1pt} de \pageref{LastPage}}

\renewcommand{\headrulewidth}{0.5pt}
\renewcommand{\footrulewidth}{0.5pt}
%%%%%%%%%%%%%%%%%%%%%%%%%%%%%%%%%%%%%%%%%%%%%%%%%%%%%%%%%%
%%%%%%%%%%%%%%%%%%%%%%%%%%%%%%%%%%%%%%%%%%%%%%%%%%%%%%%%%%

\begin{document}
    \maketitle
    \thispagestyle{fancy}

    % Problema 1
    \begin{exercise}
      Calcula las componentes del tensor de Riemann, el tensor de Ricci y el escalar de curvatura de las siguientes métrica:

      \begin{enumerate}[label = \alph*)]
        \item \(\vecspace{R}[2]\) en coordenadas cartesianas
        
        \begin{equation*}
          \odif{s^{2}} = \odif{x^{2}} + \odif{y^{2}}.
        \end{equation*}

        \begin{solution}
          Sabemos que las componentes del tensor de Riemann están dadas por

          \begin{equation}
            \tensor{R}{^{\alpha}_{\beta\mu\nu}} = \pdif{\mu}\tensor{\Gamma}{^{\alpha}_{\beta\nu}} - \pdif{\nu}\tensor{\Gamma}{^{\alpha}_{\beta\mu}} + \tensor{\Gamma}{^{\rho}_{\beta\nu}}\tensor{\Gamma}{^{\alpha}_{\rho\mu}} - \tensor{\Gamma}{^{\rho}_{\beta\mu}}\tensor{\Gamma}{^{\alpha}_{\rho\nu}}.
            \label{eq:RiemannTensorComponents}
          \end{equation}

          Para esta métrica los símbolos de Christoffel se anulan, \idest

          \begin{equation}
            \tensor{\Gamma}{^{\beta}_{\alpha\mu}} = 0.
            \label{eq:ChristoffelSymbolsForCartesianPlane}
          \end{equation}

          Sustituyendo \cref{eq:ChristoffelSymbolsForCartesianPlane} en \cref{eq:RiemannTensorComponents}, tenemos que las componentes del tensor de Riemann son:

          \begin{empheq}[box = \color{pinkwave}\fbox]{equation}
            \tensor{R}{^{\alpha}_{\beta\mu\nu}} = 0.
            \label{eq:RiemannTensorForCartesianPlane}
          \end{empheq}

          Por otro lado, sabemos que el tensor de Ricci está dado por

          \begin{equation*}
            \tensor{R}{_{\mu\nu}} = \tensor{R}{^{\sigma}_{\mu\sigma\nu}}.
          \end{equation*}

          Y puesto que el tensor de Ricci es una contracción del tensor de Riemann, \cref{eq:RiemannTensorForCartesianPlane} implica que

          \begin{empheq}[box = \color{pinkwave}\fbox]{equation*}
            \tensor{R}{_{\mu\nu}} = 0.
          \end{empheq}

          Finalmente, tenemos que el escalar de Ricci es una contracción del tensor de Ricci, por lo que

          \begin{empheq}[box = \color{pinkwave}\fbox]{equation*}
            R = 0.
          \end{empheq}
        \end{solution}
        
        \item \(\vecspace{R}[2]\) en coordenadas polares
        
        \begin{equation*}
          \odif{s^{2}} = \odif{r^{2}} + r^{2}\odif{\theta^{2}}.
        \end{equation*}

        \begin{solution}
          Notamos que la métrica describe el espacio plano en coordenadas polares y, del inciso anterior sabemos que la curvatura del espacio es 0, así
          
          \begin{empheq}[box = \color{pinkwave}\fbox]{align*}
            \tensor{R}{^{\alpha}_{\beta\mu\nu}} &= 0,\\
            \tensor{R}{_{\mu\nu}} &= 0,\\
            R &= 0.
          \end{empheq}
        \end{solution}
        
        \item La 2-esfera \(S^{2}\) en coordenadas esféricas
        
        \begin{equation*}
          \odif{s^{2}} = \odif{\theta^{2}} + \sin^{2}\theta\odif{\phi^{2}}.
        \end{equation*}

        \begin{solution}
          Para esta métrica sabemos que los símbolos de Christoffel que no se anulan son:

          \begin{align*}
            \tensor{\Gamma}{^{\theta}_{\phi\phi}} &= -\sin\theta\cos\theta,\\
            \tensor{\Gamma}{^{\phi}_{\phi\theta}} &= \tensor{\Gamma}{^{\phi}_{\theta\phi}} = \cot\theta.
          \end{align*}

          Por lo que las componentes del tensor de Riemann que debemos calcular son:

          \begin{align}
            \tensor{R}{^{\theta}_{\beta\theta\phi}} &= \pdif{\theta}\tensor{\Gamma}{^{\theta}_{\beta\phi}} + \tensor{\Gamma}{^{\rho}_{\beta\phi}}\tensor{\Gamma}{^{\theta}_{\rho\theta}} - \tensor{\Gamma}{^{\rho}_{\beta\theta}}\tensor{\Gamma}{^{\theta}_{\rho\phi}},\label{eq:ThetaRiemannComponentsFor2Sphere}\\
            \tensor{R}{^{\phi}_{\beta\phi\theta}} &= - \pdif{\theta}\tensor{\Gamma}{^{\phi}_{\beta\phi}} + \tensor{\Gamma}{^{\rho}_{\beta\theta}}\tensor{\Gamma}{^{\phi}_{\rho\phi}} - \tensor{\Gamma}{^{\rho}_{\beta\phi}}\tensor{\Gamma}{^{\phi}_{\rho\theta}}.\label{eq:PhiRiemannComponentsFor2Sphere}
          \end{align}

          De \cref{eq:ThetaRiemannComponentsFor2Sphere} tenemos

          \begin{align*}
            \tensor{R}{^{\theta}_{\phi\theta\phi}} &= \pdif{\theta}\tensor{\Gamma}{^{\theta}_{\phi\phi}} + \tensor{\Gamma}{^{\rho}_{\phi\phi}}\tensor{\Gamma}{^{\theta}_{\rho\theta}} - \tensor{\Gamma}{^{\rho}_{\phi\theta}}\tensor{\Gamma}{^{\theta}_{\rho\phi}},\\
            &= \pdif{\theta}(-\sin\theta\cos\theta) + \tensor{\Gamma}{^{\theta}_{\phi\phi}}\tensor{\Gamma}{^{\theta}_{\theta\theta}} + \tensor{\Gamma}{^{\phi}_{\phi\phi}}\tensor{\Gamma}{^{\theta}_{\phi\theta}} - \tensor{\Gamma}{^{\theta}_{\phi\theta}}\tensor{\Gamma}{^{\theta}_{\theta\phi}} - \tensor{\Gamma}{^{\phi}_{\phi\theta}}\tensor{\Gamma}{^{\theta}_{\phi\phi}},\\
            &= -\cos^{2}\theta + \cot\theta\sin\theta\cos\theta,\\
            \Acolorboxed{\tensor{R}{^{\theta}_{\phi\theta\phi}} &= \sin^{2}\theta = - \tensor{R}{^{\theta}_{\phi\phi\theta}}.}
          \end{align*}

          Y, de \cref{eq:PhiRiemannComponentsFor2Sphere},

          \begin{align*}
            \tensor{R}{^{\phi}_{\theta\phi\theta}} &= -\pdif{\theta}\tensor{\Gamma}{^{\phi}_{\theta\phi}} + \tensor{\Gamma}{^{\rho}_{\theta\theta}}\tensor{\Gamma}{^{\phi}_{\rho\phi}} - \tensor{\Gamma}{^{\rho}_{\theta\phi}}\tensor{\Gamma}{^{\phi}_{\rho\theta}},\\
            &= -\pdif{\theta}(\cot\theta) + \tensor{\Gamma}{^{\theta}_{\theta\theta}}\tensor{\Gamma}{^{\phi}_{\theta\phi}} + \tensor{\Gamma}{^{\phi}_{\theta\theta}}\tensor{\Gamma}{^{\phi}_{\phi\phi}} - \tensor{\Gamma}{^{\theta}_{\theta\phi}}\tensor{\Gamma}{^{\phi}_{\theta\theta}} - \tensor{\Gamma}{^{\phi}_{\theta\phi}}\tensor{\Gamma}{^{\phi}_{\phi\theta}},\\
            &= -\left(\dfrac{1}{\sin^{2}\theta}\right) - \cot^{2}\theta,\\
            \Acolorboxed{\tensor{R}{^{\phi}_{\theta\phi\theta}} &= 1 = - \tensor{R}{^{\phi}_{\theta\theta\phi}}.}
          \end{align*}

          De esta manera, las componentes restantes del tensor de Riemann se anulan.

          Ahora, las componentes del tensor de Ricci que debemos calcular son:

          \begin{align*}
            \tensor{R}{_{\phi\phi}} &= \tensor{R}{^{\sigma}_{\phi\sigma\phi}} = \tensor{R}{^{\theta}_{\phi\theta\phi}} + \tensor{R}{^{\phi}_{\phi\phi\phi}},\\
            \tensor{R}{_{\phi\theta}} &= \tensor{R}{_{\theta\phi}} = \tensor{R}{^{\sigma}_{\phi\sigma\theta}} = \tensor{R}{^{\theta}_{\phi\theta\theta}} + \tensor{R}{^{\phi}_{\phi\phi\theta}},\\
            \tensor{R}{_{\theta\theta}} &= \tensor{R}{^{\sigma}_{\theta\sigma\theta}} = \tensor{R}{^{\theta}_{\theta\theta\theta}} + \tensor{R}{^{\phi}_{\theta\phi\theta}}.
          \end{align*}

          Así,

          \begin{empheq}[box = \color{pinkwave}\fbox]{align*}
            \tensor{R}{_{\phi\phi}} &= \sin^{2}\theta,\\
            \tensor{R}{_{\phi\theta}} &= \tensor{R}{_{\theta\phi}} = 0,\\
            \tensor{R}{_{\theta\theta}} &= 1.
          \end{empheq}

          Finalmente, calculamos el escalar de Ricci, tal que,

          \begin{equation*}
            R = \tensor{g}{^{\mu\nu}}\tensor{R}{_{\mu\nu}} = \tensor{g}{^{\theta\theta}}\tensor{R}{_{\theta\theta}} + \tensor{g}{^{\phi\phi}}\tensor{R}{_{\phi\phi}}.
          \end{equation*}

          Así,

          \begin{empheq}[box = \color{pinkwave}\fbox]{equation*}
            R = 2.
          \end{empheq}
        \end{solution}
      \end{enumerate}
    \end{exercise}
    
    % Problema 2
    \pagebreak
    \begin{exercise}
      Partiendo de la identidad de Bianchi

      \begin{equation}
        \cdif{\lambda}\tensor{R}{^{\alpha}_{\beta\mu\nu}} 
        + \cdif{\mu}\tensor{R}{^{\alpha}_{\beta\nu\lambda}}
        + \cdif{\nu}\tensor{R}{^{\alpha}_{\beta\lambda\mu}}
        = 0
        \label{eq:BianchiIdentity}
      \end{equation}

      demuestra que el tensor de Ricci cumple

      \begin{equation*}
        \cdif[order={\mu}]{\phantom{i}}\left(\tensor{R}{_{\mu\nu}} - \dfrac{1}{2}R\tensor{g}{_{\mu\nu}}\right) = 0.
      \end{equation*}

      \emph{Hint:} contrae los índices \(\alpha\) y \(\mu\) de la identidad de Bianchi y usa que la conexión de Levi Civita es compatible con la métrica, así como las simetrías de los tensores involucrados.

      \begin{solution}
        Comenzamos contrayendo \(\alpha\) con \(\mu\) en cada uno de los términos de \cref{eq:BianchiIdentity}.

        \begin{align}
          \cdif{\lambda}\tensor{R}{^{\mu}_{\beta\mu\nu}} + \cdif{\mu}\tensor{R}{^{\mu}_{\beta\nu\lambda}} + \cdif{\nu}\tensor{R}{^{\mu}_{\beta\lambda\mu}} &= 0,\nonumber\\
          \cdif{\lambda}\tensor{R}{_{\beta\nu}} + \cdif{\alpha}\tensor{R}{^{\alpha}_{\beta\nu\lambda}} - \cdif{\nu}\tensor{R}{_{\beta\lambda}} &= 0,\label{eq:FirstContractionInBianchiIdentity}
        \end{align}

        donde \(\tensor{R}{_{\beta\nu}}\) es el tensor de Ricci y \(\tensor{R}{_{\beta\lambda}}\) es el negativo de este, por la contracción 3) que se encuentra en las notas.

        Usando la compatibilidad de la métrica, \idest \(\cdif{\alpha} \tensor{g}{^{\beta\lambda}} = 0\) podemos contraer más índices. Así, \cref{eq:FirstContractionInBianchiIdentity} queda como:

        \begin{align}
          \tensor{g}{^{\beta\lambda}}\cdif{\lambda}\tensor{R}{_{\beta\nu}} + \tensor{g}{^{\beta\lambda}}\cdif{\alpha}\tensor{R}{^{\alpha}_{\beta\nu\lambda}} - \tensor{g}{^{\beta\lambda}}\cdif{\nu}\tensor{R}{_{\beta\lambda}} &= 0,\nonumber\\
          \cdif{\lambda}\left(\tensor{g}{^{\beta\lambda}}\tensor{R}{_{\beta\nu}}\right) + \cdif{\alpha}\left(\tensor{g}{^{\beta\lambda}}\tensor{R}{^{\alpha}_{\beta\nu\lambda}}\right) - \cdif{\nu}\left(\tensor{g}{^{\beta\lambda}}\tensor{R}{_{\beta\lambda}}\right) &= 0,\nonumber\\
          \cdif{\lambda}\tensor{R}{^{\lambda}_{\nu}} + \cdif{\alpha}\tensor{R}{^{\alpha}_{\nu}} - \cdif{\nu}R &= 0.\label{eq:SecondContractionInBianchiIdentity} 
        \end{align}

        Escribiendo \(\cdif{\lambda}\) como
        
        \begin{equation}
          \cdif{\lambda} = \cdif[order = {\mu}]{\phantom{i}}\tensor{g}{_{\lambda\mu}}.
          \label{eq:LoweringIndexInLeviCivitaConection}
        \end{equation}

        \pagebreak
        Por \cref{eq:LoweringIndexInLeviCivitaConection}, \cref{eq:SecondContractionInBianchiIdentity} queda

        \begin{align*}
          \cdif[order = {\mu}]{\phantom{i}}\tensor{g}{_{\lambda\mu}}\tensor{R}{^{\lambda}_{\nu}} + \cdif[order = {\mu}]{\phantom{i}}\tensor{g}{_{\alpha\mu}}\tensor{R}{^{\alpha}_{\nu}} - \cdif[order = {\mu}]{\phantom{i}}\tensor{g}{_{\nu\mu}}R &= 0,\\
          \cdif[order = {\mu}]{\phantom{i}}\tensor{R}{_{\mu\nu}} + \cdif[order = {\mu}]{\phantom{i}}\tensor{R}{_{\mu\nu}} - \cdif[order = {\mu}]{\phantom{i}}\tensor{g}{_{\mu\nu}}R &= 0,\\
          2\cdif[order = {\mu}]{\phantom{i}}\tensor{R}{_{\mu\nu}} - \cdif[order = {\mu}]{\phantom{i}}R\tensor{g}{_{\mu\nu}} &= 0,\\
          \cdif[order = {\mu}]{\phantom{i}}\left(2\tensor{R}{_{\mu\nu}} - R\tensor{g}{_{\mu\nu}}\right) &= 0,\\
          \Acolorboxed{\cdif[order = {\mu}]{\phantom{i}}\left(\tensor{R}{_{\mu\nu}} - \dfrac{1}{2}R\tensor{g}{_{\mu\nu}}\right) &= 0.}
        \end{align*}
      \end{solution}
    \end{exercise}
    
    % Problema 3
    \pagebreak
    \begin{exercise}
      El ansatz para la solución	 de Schwarzschild (esféricamente simétrica) que discutimos en clase se escribe en coordenadas polares como

      \begin{equation}
        \odif{s^{2}} = -f(r)\odif{t^{2}} + h(r)\odif{r^{2}} + r^{2}(\odif{\theta^{2}} + \sin^{2}\theta\odif{\varphi^{2}}).
        \label{eq:MetricProblem3}
      \end{equation}

      Sustituir esta métrica en las ecuaciones de Einstein lleva a las siguientes tres ecuaciones diferenciales independientes (las demás componentes se satisfacen automáticamente)

      \begin{align}
        -rh(r)\tensor{f}{^{\prime}}(r)^{2} + f(r)(2h(r)(r\tensor{f}{^{\prime\prime}}(r) + \tensor{f}{^{\prime}}(r)) - r\tensor{f}{^{\prime}}(r)\tensor{h}{^{\prime}}(r)) - 2f(r)^{2}\tensor{h}{^{\prime}}(r) &= 0,\label{eq:DiffEq1}\\
        r\tensor{h}{^{\prime}}(r) + h(r)^{2} - h(r) &= 0,\label{eq:DiffEq2}\\
        r\tensor{f}{^{\prime}}(r) - f(r)h(r) + f(r) &= 0.\label{eq:DiffEq3}
      \end{align}

      En este problema resolverás este sistema de ecuaciones.

      \begin{enumerate}[label = \alph*)]
        \item Usa (\ref{eq:DiffEq2}) y (\ref{eq:DiffEq3}) para demostrar que
        
        \begin{equation}
          f(r) = \dfrac{k}{h(r)},
        \end{equation}

        con \(k\) una constante arbitraria que fijaremos a 1 sin pérdida de generalidad.

        \begin{solution}
          Primero, multiplicamos \cref{eq:DiffEq2} por \(f(r)\) y \cref{eq:DiffEq3} por \(h(r)\).

          \begin{align}
            rh^{\prime}f + fh^{2} - fh &= 0,\label{eq:DiffEq2ByF}\\
            rf^{\prime}h + fh^{2} + fh &= 0,\label{eq:DiffEq3ByH}
          \end{align}

          donde \(f(r) = f\) y \(h(r) = h\).

          Ahora, sumamos \cref{eq:DiffEq2ByF,eq:DiffEq3ByH},

          \begin{align}
            r(h^{\prime}f + hf^{\prime}) &= 0,\nonumber\\
            r(hf)^{\prime} &= 0,\nonumber\\
            \odv[fun=true]{hf}{r} &= 0.\label{eq:DiffEqOfSum}
          \end{align}

          Integrando \cref{eq:DiffEqOfSum} respecto a \(r\).

          \begin{align*}
            h(r)f(r) + C &= 0,\\
            h(r)f(r) &= k,\\
            f(r) &= \dfrac{k}{h(r)}.
          \end{align*}

          Igualando \(k = 1\), la expresión anterior queda como:

          \begin{empheq}[box = \color{pinkwave}\fbox]{equation}
            f(r) = \dfrac{1}{h(r)}.
            \label{eq:FInTermsOfH}
          \end{empheq}
        \end{solution}

        \item Esto te deja con 2 ecuaciones para una sola variable \(h(r)\), una de primer orden y otra de segundo orden. Demuestra que no son independientes derivando la de primer orden respecto a \(r\) y sustituyendo en la de segundo orden.
        
        \begin{solution}
          Desarrollando \cref{eq:DiffEq1} y usando que \(hf = 1\) obtenemos que:

          \begin{align}
            -rh{f^{\prime}}^{2} + 2rf^{\prime\prime}fh + 2f^{\prime}fh - rf^{\prime}fh^{\prime} - 2f^{2}h^{\prime} &= 0,\nonumber\\
            -rh{f^{\prime}}^{2} + 2rf^{\prime\prime} + 2f^{\prime} -rf^{\prime}fh^{\prime} - 2f^{2}h^{\prime} &= 0.\label{eq:DiffEq2Reduce1}
          \end{align}

          Calculamos la primera y segunda derivada de \cref{eq:FInTermsOfH},

          \begin{align*}
            f^{\prime} &= \dfrac{-h^{\prime}}{h^{2}},\\
            f^{\prime\prime} &= \dfrac{-h^{\prime\prime}h^{2} + 2{h^{\prime}}^{2}h}{h^{4}}.
          \end{align*}

          Sustituyendo las derivadas correspondientes en \cref{eq:DiffEq2Reduce1},

          \begin{align}
            -rh\dfrac{{h^{\prime}}^{2}}{h^{4}} + 2r\left(\dfrac{-h^{\prime\prime}h^{2} + 2{h^{\prime}}^{2}h}{h^{4}}\right) - r\left(\dfrac{-h^{\prime}}{h^{2}}\right)\left(\dfrac{1}{h}\right)h^{\prime} - 2\left(\dfrac{1}{h^{2}}\right)h^{\prime} &= 0,\nonumber\\
            \dfrac{rh^{\prime\prime}h^{2} - 2r{h^{\prime}}^{2}h + 2h^{\prime}h^{2}}{h^{4}} &= 0,\nonumber\\
            \Aboxed{h(r)(rh^{\prime\prime}(r) + 2h^{\prime}(r)) - 2rh^{\prime}(r)^{2} &= 0.}\label{eq:DiffEq4}
          \end{align}

          Derivando \cref{eq:DiffEq2} respecto a \(r\),

          \begin{align*}
            \odv[fun=true]{rh^{\prime} + h^{2} - h}{r} &= h^{\prime} + rh^{\prime\prime} + 2hh^{\prime} - h^{\prime},\\
            &= rh^{\prime\prime} + 2h(r)h^{\prime}.
          \end{align*}

          Puesto que la expresión anterior aparece en \cref{eq:DiffEq4}, podemos concluir que no son independientes, ya que \cref{eq:DiffEq4} se puede escribir como combinación lineal de \cref{eq:DiffEq2}.
        \end{solution}
        
        \item Finalmente demuestra que la solución general a la única ecuación de primer orden para \(h(r)\) es
        
        \begin{equation}
          h(r) = \left(1 + \dfrac{C}{r}\right)^{-1},
        \end{equation}

        donde \(C\) es una constante de integración arbitraria.

        \begin{solution}
          Finalmente, resolvemos \cref{eq:DiffEq3} para encontrar la solución general. Primero, reescribimos la ec. diferencial, tal que,

          \begin{equation*}
            rh^{\prime} - h = -h^{2}.
          \end{equation*}

          Notamos que es una ec. diferencial de Bernoulli. Así, \(u = h^{-1}\ \implies h = u^{-1}\ \wedge h^{\prime} = -u^{\prime}u^{-2}\). Sustituyendo

          \begin{equation*}
            u^{\prime} + \dfrac{1}{r}u = \dfrac{1}{r}.
          \end{equation*}

          Ahora, para resolver esta nueva ecuación diferencial, tenemos

          \begin{align*}
            \mu &= \e^{\int\frac{1}{r}\odif{r}},\\
            &= \e^{\log(r) + C_{1}},\\
            &= \e^{\log(r)}\e^{C_{1}},\\
            \Aboxed{\mu &= rC_{2}.}
          \end{align*}

          \pagebreak
          Así,

          \begin{align*}
            u &= \int \dfrac{1}{r}rC_{2}\odif{r},\\
            &= C_{2}\int\odif{r},\\
            \Aboxed{u &= \dfrac{r + C}{r}.}
          \end{align*}

          Y sabemos que \(h^{-1} = u\), entonces

          \begin{empheq}[box = \fbox]{equation*}
            h^{-1}(r) = 1 + \dfrac{C}{r}.
          \end{empheq}

          Por lo tanto, la solución general es

          \begin{empheq}[box = \color{pinkwave}\fbox]{equation*}
            h(r) = \left(1 + \dfrac{C}{r}\right)^{-1}.
          \end{empheq}
        \end{solution}
      \end{enumerate}
    \end{exercise}
\end{document}